\documentclass[11pt]{article}
\usepackage[a4paper,margin=28mm]{geometry}
\usepackage{amsmath,amssymb,amsthm,hyperref}
\newtheorem{theorem}{Theorem}
\newtheorem{lemma}{Lemma}
\newcommand{\fall}[2]{#1^{\underline{#2}}}
\title{A finite pointwise bound for alternating-block dice}
\author{Research notes, 30 September 2026}
\date{}
\begin{document}
\maketitle

\paragraph{Status and attribution.}
The construction below is a dice realization of the classical shelf shuffle,
up to reversing labels or inverting the permutation. Its pointwise mixing
scale is already established by Diaconis, Fulman, and Holmes,
\emph{Analysis of casino shelf shuffling machines}, Annals of Applied
Probability 23 (2013), 1692--1720,
\href{https://arxiv.org/abs/1107.2961}{arXiv:1107.2961}, Theorem 3.4.
The purpose here is a self-contained finite-parameter estimate suitable for
the fair-dice article. It improves that article's stated approximate
construction; it does not solve the exact-fairness open problems.

For $n\geq2$, let
\[
 B_n=12\cdots n\,n\cdots21,\qquad s_{n,m}=B_n^m.
\]
Each die has $2m$ faces. Let $P_{n,m}(\pi)$ be its order distribution.

\begin{theorem}
For $n\geq3$ and $m>n/2$, every permutation satisfies
\[
 1-\frac{n(n-1)(n-2)}{6m^2}
 \leq n!P_{n,m}(\pi)
 \leq
 \exp\left(\frac{n(n-1)(n-2)}{12(m^2-n^2/4)}\right).
\]
Consequently, for $0<\varepsilon\leq1$, the choice
\[
 m=\left\lceil\sqrt{\frac{n^2}{4}
          +\frac{n(n-1)(n-2)}{6\varepsilon}}\right\rceil
\]
gives $|n!P_{n,m}(\pi)-1|\leq\varepsilon$ for every $\pi$.
Thus $O(n^{3/2}\varepsilon^{-1/2})$ faces per die suffice.
For $n=2$, the word $1221$ is already exactly fair.
\end{theorem}

\begin{proof}
Represent each chosen face by a uniform level in $[m]$ and a uniform choice
of one of the two occurrences within that level. A block restricted to any
two letters is pairwise fair. Conditional on every occupied level containing
at most two dice, the output permutation is uniform: the allocation of dice
to levels is invariant under relabelling, and the order of each two-die cell
is a fair independent binary choice. A union bound over triples therefore
gives
\[
 P_{n,m}(\pi)\geq\frac1{n!}
 \Pr(\text{no level contains three dice})
 \geq\frac1{n!}\left(1-\binom n3m^{-2}\right).
\]

For the upper bound, a prescribed order of $r\geq1$ different letters has
at most two occurrences as a subsequence of $B_n$. Indeed, such an order
must increase and then decrease, and the only possible ambiguity is which
of the two occurrences of its largest letter is selected. Its conditional
probability within one level is therefore at most $2^{1-r}$.
For any fixed full permutation, distributing its consecutive segments among
the $m$ levels gives
\[
 P_{n,m}(\pi)
 \leq m^{-n}[z^n]
       \left(1+\sum_{r\geq1}2^{1-r}z^r\right)^m
 =m^{-n}[z^n]\left(\frac{1+z/2}{1-z/2}\right)^m.
\]
In fact equality holds for the increasing permutation, because every
consecutive segment of that permutation is admissible.
After the substitution $t=mz$, the right side, multiplied by $n!$, is
\[
 n![t^n]\exp\left(t+
       \sum_{j\geq1}\frac{t^{2j+1}}
                 {(2j+1)2^{2j}m^{2j}}\right).
\]
When the exponential is expanded, a choice of terms of total degree $d$
contributes the falling factorial $\fall nd$. The elementary inequality
\[
 \fall n{d_1+\cdots+d_a}
 \leq\prod_{b=1}^a\fall n{d_b}
 \qquad(d_b\geq0,\ \sum_b d_b\leq n)
\]
allows the contributions of distinct non-linear terms to be separated.
Terms whose total degree exceeds $n$ only increase the resulting bound.
It follows that
\[
 n!P_{n,m}(\pi)\leq
 \exp\left(\sum_{1\leq j\leq(n-1)/2}
 \frac{\fall n{2j+1}}{(2j+1)2^{2j}m^{2j}}\right).
\]
Since $\fall n{2j+1}\leq\fall n3\,n^{2j-2}$ and $2j+1\geq3$, this is at
most
\[
 \exp\left(\frac{\fall n3}{12m^2}
       \sum_{j\geq1}\left(\frac{n^2}{4m^2}\right)^{j-1}\right),
\]
which is the asserted upper bound. For the specified $m$, its exponent is
at most $\varepsilon/2$, and $e^{\varepsilon/2}\leq1+\varepsilon$ on
$0\leq\varepsilon\leq1$. The lower error is at most $\varepsilon$.
\end{proof}

\paragraph{A finite obstruction for this family.}
The coefficient calculation for the increasing permutation is exact and
has nonnegative terms. Keeping the first nonconstant term gives
\[
 n!P_{n,m}(12\cdots n)-1
 \geq\frac{n(n-1)(n-2)}{12m^2}.
\]
Thus relative error at most $\varepsilon$ in this family requires
$m^2\geq n(n-1)(n-2)/(12\varepsilon)$. The sufficient bound above is
within an asymptotic factor $\sqrt2$ in the number of faces when
$n/\varepsilon\to\infty$. This also shows directly, without an asymptotic
shuffle theorem, the sharp $n^{3/2}/\sqrt\varepsilon$ scale for this family.

\paragraph{Verification.}
The accompanying \texttt{verify\_shelf.py} computes subsequence counts using
integer arithmetic, checks every permutation for small $n,m$, compares the
maximum with the generating-function formula, and checks the stated bounds.
The proof, rather than numerical agreement, establishes the arbitrary-size
claims.
\end{document}
